% \section{Algorithm}
% \label{Section: Methodology}
% \subsection{DualHash: A Stochastic Primal-Dual Hashing Algorithm}
% Inspired by PDHG \cite{chambolleFirstOrderPrimalDualAlgorithm2011} in convex optimization and its nonconvex extension PPDG \cite{guoPreconditionedPrimalDualGradient2023}, we consider the equivalent constrained formulation of  the problem \eqref{Prop: block-wise regularized problem}:
% \begin{align} 
% \min_{\bx \in \mathcal{X}, \bB \in \mathcal{B}, \by \in \mathcal{Y}} &\quad F(\bx, \bB) + h(\bY) \notag \\
% \text{s.t. } &\quad \bY = \bB \nonumber ,
% \end{align} whose corresponding dual problem is 
% \begin{align}
% \label{Prop: the dual problem}
% \max_{\bLmbd \in \mathcal{Y}^*} \min_{\bx \in \mathcal{X}, \bB \in \mathcal{B}, \by \in \mathcal{Y}} &\quad F(\bx, \bB) + h(\bY) + \langle \bLmbd, \bB - \bY \rangle. \tag{D}
% \end{align} 
% Recalling the definition of the conjugate function \cite{Rockafellar1998VariationalA}:
% \begin{equation}
% \label{eq:fenchel_conjugate}
% h^*(\bY) := \sup_{\bx \in \mathcal{Y}} \{ \langle \bY, \bx \rangle - h(\bx) \}, \quad \bY \in \mathcal{Y}^*
% \end{equation} and observing that the minimization with respect to $\bY$ in the problem \eqref{Prop: the dual problem} can be performed independently, we then simplify the Lagrangian function in the problem \eqref{Prop: the dual problem} to 
% \begin{align}
% \label{Eq: simplified_lagrangian}
% \calL(\bx, \bB, \bLmbd) := F(\bx, \bB) + \langle \bB, \bLmbd \rangle - h^*(\bLmbd).\tag{\(\calL\)}
% \end{align}

% Through this transformation, we replace $h(\bY)$ with $-h^*(\bLmbd)$ and can optimize \(h^*\) in the dual space. Since $h^*$ is always convex and l.s.c. \cite{Rockafellar1998VariationalA}, it means that the $\text{prox}_{h^*}(\cdot)$ is often much easier to compute (potentially having a closed-form solution) than that of $h$.  For $h(\bz) = \sum_{i=1}^n \lambda ||z_i| - 1|$, we derive its conjugate function as illustrated in \Cref{fig:w_type_regularization}:
% \begin{align*}
% &h^*(\bx) = \sum_{i=1}^n h^*(x_i),\,\, \notag \\
% & \text{where } h^*(x_i) = \begin{cases}
%     |x_i|, & \text{if } x_i \in [-\lambda, \lambda],\\
%     +\infty, & \text{otherwise},
%     \end{cases}
% \end{align*} whose proximal operator  in the piecewise solution is 
% \begin{align}
%     \label{Eq: the proximal mapping of conjugate function}
%     [\operatorname{prox}_{\tau h^*}(\by)]_i = 
%     \begin{cases}
%         \lambda, & y_i > \lambda +\tau,\\
%         y_i -\tau, & \tau < y_i \leq \lambda +\tau, \\
%         0, & -\tau \leq y_i \leq\tau, \\
%         y_i +\tau, &  - \lambda -\tau \leq y_i < -\tau, \\
%         -\lambda, & y_i < -\lambda -\tau.
%     \end{cases}
% \end{align} This motivates us to perform a proximal step to update \(\bLmbd\).
% \begin{figure}[!t]
% \centering
% \includegraphics[width=\columnwidth]{papers_img/W_type_with_its_conjugate_func.pdf}
% \caption{W-type regularization function $h(z) = \lambda ||z| - 1|$ and its conjugate function $h^*(x)$.}
% \label{fig:w_type_regularization}
% \vspace{-1ex}
% \end{figure} We now outline \textbf{DualHash}, which employs an alternating update scheme for $\bx, \bB, \bm{\Lambda}$. 
% we adopt a stochastic first-order method for the $\bx$-update since computing the full gradient \( \nabla_{\bx} F(\bx, \bB) \) is computationally expensive for typically the large number of samples $n$ in deep learning. Specifically, at each iteration $k$, we uniformly sample a mini-batch \(\calJ_k \subseteq \{1,\cdots,n\}\) with \(b_k = |\calJ_k| \leq n\) and compute:
% \begin{align}
% \label{Eq: the mini batch stochastic gradient estimator of F at x}
% \nabla_{\bx} F(\bx^k, \bB^k; \calJ_k)
% =\, \frac{1}{|\mathcal{J}_k|}\sum_{j \in \mathcal{J}_k} \nabla_{\bx} F_j(\bx^k, \bb_j^k).
% \end{align}
% Then we utilize this approximation within an optimization strategy chosen to enhance efficiency (e.g., incorporating momentum or variance reduction), as detailed in \Cref{Sec: Algorithm implementations}.
% Second, the auxiliary variable $\bB$ is updated via a simple gradient descent step \eqref{Eq: the update of B}; this suffices because the objective \eqref{Eq: simplified_lagrangian} with respect to $\bB$ is smooth with an inexpensive gradient, 
% % which is slightly different from \cite{chambolleFirstOrderPrimalDualAlgorithm2011}:
%     \begin{align}
%    \nonumber \bB^{k+1} &= \bB^k - \tau \nabla_{\bB}\calL(\bx^{k+1}, \bB^k, \bLmbd^k)\\
%     &= \bB^k - \tau \left( \nabla_{\bB} F(\bx^{k+1}, \bB^k) + \bLmbd^k \right). \label{Eq: the update of B}
%     \end{align}
%     Finally, the dual variable $\bLmbd$ is updated via proximal gradient ascent on $h^*$, using an extrapolation step introduced in \cite{chambolleFirstOrderPrimalDualAlgorithm2011}. The update solves the following optimization problem:
% \begin{align}
%     \bLmbd^{k+1} &= \text{prox}_{(1/\tau h^*)} \left( \bLmbd^k + \left(2\bB^{k+1} - \bB^k  \right) \right) \notag \\
%     &=  \arg\min_{\bLmbd \in \mathcal{Y}^* } \left\{ h^*(\bLmbd) - \langle \bLmbd, 2\bB^{k+1} - \bB^k \rangle \right. \notag \\
%     &\quad\quad\quad\quad\quad\quad\quad\quad\quad\quad\left. + \frac{\tau}{2} \| \bLmbd - \bLmbd^k\|^2 \right\}. 
%     \label{Eq: the update of lambda}
% \end{align}
% \subsection{Implementations}
% \label{Sec: Algorithm implementations}
% The unified DualHash framework allows for flexible choices in the $\bx$ update step based on a mini-batch stochastic gradient approximation \eqref{Eq: the mini batch stochastic gradient estimator of F at x}. This flexibility enables us to incorporate advanced stochastic optimization techniques. We thus instantiate the $\bx$-update using two well-established approaches: \textbf{momentum} \cite{Nesterov1983AMF, Polyak1964} and \textbf{variance reduction} \cite{driggs2022biased}, which can improve empirical training efficiency and can also improve the theoretical convergence guarantees \cite{Liu2020AnIA,sutskever2013importance}.

% % \vspace{5pt}
% \paragraph{DualHash-StoM}
% %Numerically, 
% The heavy ball method \cite{Polyak1964} and Nesterov acceleration \cite{Nesterov1983AMF} can be integrated with stochastic gradient descent (SGD) to yield accelerated variants like SGD with momentum, which significantly speeds up convergence.
% This approach has become a popular strategy for training neural networks \cite{fotopoulos2024review,Hertrich2020InertialSP}. We thus incorporate this technique and compute the stochastic gradient estimator (Line 3 of \Cref{alg:our_algorithm_sgdm}) as:
% \begin{align}
% \label{Eq: the stochastic approximate in SGDM}
%     \bG^k = \nabla_{\bx} F(\bz^k, \bB^k; \calJ_k),\, \text{where } \bz^k = \bx^k + \beta_k (\bx^k - \bx^{k-1}).
% \end{align} The update of \( \bx^k\) (Line 4 of Alg. \ref{alg:our_algorithm_sgdm}) is then:
% \begin{align}
% \label{Eq: the update of x in SGDM}
%     \bx^{k+1} = \by^k - \eta_k \bG^k,\,\text{where } \by^k = \bx^k + \alpha_k (\bx^k - \bx^{k-1}). 
% \end{align} The overall algorithm, DualHash-StoM, is detailed in \Cref{alg:our_algorithm_sgdm}.
% \begin{remark}
% The momentum acceleration framework in \eqref{Eq: the update of x in SGDM} covers these two accelerations:
% Nesterov acceleration \cite{Nesterov1983AMF} corresponds to $\alpha_k = \beta_k \neq 0$; the heavy-ball method \cite{Polyak1964} is obtained with $\beta_k = 0$; vanilla SGD is recovered when $\alpha_k = \beta_k = 0$; In this paper, we set $\alpha_k, \beta_k \in (0,1)$ for all $k \geq 1$. 
% \end{remark}

% % \vspace{3pt}
% \paragraph{DualHash-StoRM}
% Alternatively, we employ the STORM estimator \cite{cutkoskyMomentumBasedVarianceReduction2019} for the $\bx$-update. STORM utilizes a recursive momentum mechanism to reduce gradient variance, which is motivated by SGDM. Unlike some other variance reduction techniques (e.g., SVRG \cite{johnsonAcceleratingStochasticGradient2013}, SARAH \cite{SARAH2017}) that require full-gradient or large-batch gradient computations, STORM achieves variance reduction without such requirements. Given its practical effectiveness in deep learning \cite{cutkoskyMomentumBasedVarianceReduction2019,levy2021storm+}, we compute the stochastic gradient estimator (Line 3 of \Cref{alg:our_algorithm_storm}) as:
% % {\small
% % \begin{align}
% % \label{Eq: the stochastic approximate in STORM}
% % \bD^k =
% % \begin{cases} 
% % \nabla_{\bx} F(\bx^1, \bB^1; \calJ_1), & \text{if } k = 1, \\
% % \nabla_{\bx} F(\bx^k, \bB^k; \calJ_k) + (1 - \rho_{k-1}) \left( \bD^{k-1} - \nabla_{\bx} F(\bx^{k-1}, \bB^{k-1};\calJ_{k-1})\right), & \text{if }k > 1.
% % \end{cases}
% % \end{align}}
% {\small
% \begin{align}
% \label{Eq: the stochastic approximate in STORM}
% \bD^k =
% \begin{cases} 
% \nabla_{\bx} F(\bx^1, \bB^1; \mathcal{J}_1), &\text{if }  k = 1, \\
% (1 - \rho_{k-1}) \big( \bD^{k-1} - \nabla_{\bx} F(\bx^{k-1}, \bB^{k-1};\mathcal{J}_{k-1})\big)\\
% +\nabla_{\bx} F(\bx^k, \bB^k; \mathcal{J}_k), & \text{if } k > 1.
% \end{cases}
% \end{align}}
% The update of $\bx$ (Line 4 of \Cref{alg:our_algorithm_storm}) is then 
% \begin{align}
% \label{Eq: the update of x in STORM}
%     \bx^{k+1} = \bx^k - \eta_k \bD^k.
% \end{align} The overall algorithm, DualHash-StoRM, is detailed in \Cref{alg:our_algorithm_storm}.

% \begin{figure}[!htbp]
%     \vspace{-5mm} 
%     \begin{minipage}[t]{0.485\textwidth}
%         \begin{algorithm}[H]
%        \caption{DualHash-StoM}
%         \small   % 或使用 \footnotesize 更小
%         \setlength{\baselineskip}{0.8\baselineskip}  % 减少行间
%       \label{alg:our_algorithm_sgdm}
%       \begin{algorithmic}[1]
%         \STATE \textbf{Initialize:} $\bx^1 = \bz^1 = \by^1 \in \calX $, \( \bB^1 \in \calB,\, \bLmbd^1\in \calY^*\), parameters $\{\eta_k\} $, $\tau$, $b_k$, $\{\alpha_k, \beta_k \in (0,1)\}$
%         \FOR{$k = 1, 2, \ldots$}               
%           \STATE Compute \(\bG^k\) by \eqref{Eq: the stochastic approximate in SGDM}
%           \STATE Update $\bx^{k+1}$ by \eqref{Eq: the update of x in SGDM}
%           \STATE Update $\bB^{k+1}$ by \eqref{Eq: the update of B}
%           \STATE Update $\bm{\Lambda}^{k+1}$ by \eqref{Eq: the update of lambda}
%         \ENDFOR
%         \STATE \textbf{Return:} $\bx_R$ where $R$ is uniformly sampled from $\{2, \ldots, T+1\}$
%       \end{algorithmic}
%     \end{algorithm}
%   \end{minipage}
%   \hspace{0.5mm} 
%   \begin{minipage}[t]{0.485\textwidth}
%     \begin{algorithm}[H]
%         \caption{ DualHash-StoRM}
%       \small
%       \setlength{\baselineskip}{0.8\baselineskip}
%       \label{alg:our_algorithm_storm}
%       \begin{algorithmic}[1]
%         \STATE  \textbf{Initialize:} $\bx^1  \in \calX$, \( \bB^1 \in \calB,\, \bLmbd^1\in \calY^*\), parameters $\{\eta_k\} $, $\tau $, $b_k$, $\{\rho_k\in (0,1)\}$
%         \FOR{$k = 1, 2, \ldots$}
%         \STATE Compute $\bD^k$ by \eqref{Eq: the stochastic approximate in STORM}
%         \STATE Update $\bx^{k+1}$ by \eqref{Eq: the update of x in STORM}.
%         \STATE Update $\bB^{k+1}$ by \eqref{Eq: the update of B}.
%         \STATE Update $\bm{\Lambda}^{k+1}$ by \eqref{Eq: the update of lambda}.
%         \ENDFOR
%         \STATE \textbf{Return:} $\bx_R$ where $R$ is uniformly sampled from $\{2, \ldots, T+1\}$
%       \end{algorithmic}
%     \end{algorithm}
%   \end{minipage}
%   \vspace{-1ex}
% \end{figure}
% \begin{remark}
% The update steps for $\bB$ and $\bm{\Lambda}$ imply the following optimality conditions, which may result in explicit closed-solution in deep hashing. The $\bB$-update in \eqref{Eq: the update of B} yields:
% \begin{align}
%     \label{Eq: the optimality condition of Bk}
%     \bB^{k} = \bB^{k-1} - \tau \left( \nabla_{\bB}F(\bx^{k}, \bB^{k-1}) + \bm{\Lambda}^{k-1} \right).
% \end{align}
% For the $\bm{\Lambda}$-update in \eqref{Eq: the update of lambda}, the proximal operator definition gives:
% \[
%     \tau \left( \bm{\Lambda}^{k} + \tau^{-1}(2\bB^{k} - \bB^{k-1}) - \bm{\Lambda}^{k} \right) \in \partial h^{*}(\bm{\Lambda}^{k}),
% \]
% which is equivalent to the existence of $\mathcal{G}_{\bm{\Lambda}}^{k} \in \partial h^{*}(\bm{\Lambda}^{k})$ such that:
% \begin{align}
%     \label{Eq: the optimality condition of lbda}
%     \mathcal{G}_{\bm{\Lambda}}^{k} = (2\bB^{k} - \bB^{k-1}) + \tau(\bm{\Lambda}^{k-1} - \bm{\Lambda}^{k}).
% \end{align}
% If the sequence \(\{ (\bx^k, \bB^k, \bm{\Lambda}^k) \}\) converges to \((\bx^*, \bB^*, \bm{\Lambda}^*)\), then by the outer semicontinuity of $\partial h^{*}(\cdot)$, we have $\bB^* \in \partial h^*(\bm{\Lambda}^*)$.
% % which corresponds to the necessary condition \eqref{Eq: the sufficient and necessary condition for dual} for the second relation in optimality condition \eqref{Eq: the sufficient and necessary condition}.
% This establishes \eqref{Eq: the sufficient and necessary condition for dual}, which is precisely the necessary condition for the second relation in \eqref{Eq: the sufficient and necessary condition}.
% \end{remark}
\section{Algorithm}
\label{Section: Methodology}
\subsection{DualHash: A Stochastic Primal-Dual Hashing Framework}
Inspired by PDHG \cite{chambolleFirstOrderPrimalDualAlgorithm2011},
% and its nonconvex extension PPDG \cite{guoPreconditionedPrimalDualGradient2023}, 
we consider the equivalent constrained formulation of  the problem \eqref{Prop: block-wise regularized problem}:
\begin{align} 
\min_{\bx \in \mathcal{X}, \bB \in \mathcal{B}, \bY \in \mathcal{Y}} &\quad F(\bx, \bB) + h(\bY) \notag \\
\text{s.t. } &\quad \bY = \bB \nonumber ,
\end{align} whose corresponding dual problem is 
\begin{align}
\label{Prop: the dual problem}
\max_{\bLmbd \in \mathcal{Y}^*} \min_{\bx \in \mathcal{X}, \bB \in \mathcal{B}, \bY \in \mathcal{Y}} &\quad F(\bx, \bB) + h(\bY) + \langle \bLmbd, \bB - \bY \rangle. \tag{D}
\end{align} 
Recalling the definition of the conjugate function \cite{Rockafellar1998VariationalA}:
\begin{equation}
\label{eq:fenchel_conjugate}
h^*(\bY) := \sup_{\by \in \mathcal{Y}} \{ \langle \bY, \by \rangle - h(\by) \}, \quad \bY \in \mathcal{Y}^*
\end{equation} and observing that the minimization with respect to $\bY$ in the problem \eqref{Prop: the dual problem} can be performed independently, we then simplify the Lagrangian function in the problem \eqref{Prop: the dual problem} to 
\begin{equation}
\label{Eq: simplified_lagrangian}
\calL(\bx, \bB, \bLmbd) := F(\bx, \bB) + \langle \bB, \bLmbd \rangle - h^*(\bLmbd).
\end{equation}
Through this transformation, we replace $h(\bY)$ with $-h^*(\bLmbd)$ and optimize the \(h^*\) in the dual space. The convexity of $h^*$ makes $\operatorname{prox}_{h^*}(\cdot)$ often much easier to compute (potentially having a closed-form solution) than $\operatorname{prox}_{h}(\cdot)$.
For $h(\bz) = \sum_{i=1}^n \lambda ||z_i| - 1|$, we derive its conjugate function as illustrated in \Cref{fig:w_type_regularization}:
\begin{align*}
&h^*(\bx) = \sum_{i=1}^n h^*(x_i),\,\, \text{where } h^*(x_i) = \begin{cases}
    |x_i|, & \text{if } x_i \in [-\lambda, \lambda],\\
    +\infty, & \text{otherwise},
    \end{cases}
\end{align*} whose proximal operator in the piecewise solution is 
\begin{align}
    \label{Eq: the proximal mapping of conjugate function}
    [\operatorname{prox}_{\tau h^*}(\by)]_i = 
    \begin{cases}
        \lambda, & y_i > \lambda +\tau,\\
        y_i -\tau, & \tau < y_i \leq \lambda +\tau, \\
        0, & -\tau \leq y_i \leq\tau, \\
        y_i +\tau, &  - \lambda -\tau \leq y_i < -\tau, \\
        -\lambda, & y_i < -\lambda -\tau.
    \end{cases}
\end{align} This motivates us to perform a proximal step to update \(\bLmbd\).
\begin{figure}[t]
\centering
\includegraphics[width=0.6\textwidth]{papers_img/W_type_with_its_conjugate_func.pdf}
\caption{W-type regularization function $h(z) = \lambda ||z| - 1|$ and its conjugate function $h^*(x)$.}
\label{fig:w_type_regularization}
\vspace{-3 ex}
\end{figure}

We now outline the \textbf{DualHash} framework, which employs an alternating update scheme for $\bx, \bB, \bm{\Lambda}$. We adopt a stochastic first-order method for the $\bx$-update since computing the full gradient \( \nabla_{\bx} F(\bx, \bB) \) is computationally expensive for a typically large number of samples $n$ in deep learning. Specifically, at each iteration $k$, we uniformly sample a mini-batch \(\calJ_k \subseteq \{1,\cdots,n\}\) with \(b_k = |\calJ_k| \leq n\) and compute:
\begin{align}
\label{Eq: the mini batch stochastic gradient estimator of F at x}
\nabla_{\bx} F(\bx^k, \bB^k; \calJ_k)
=\, \frac{1}{|\mathcal{J}_k|}\sum_{j \in \mathcal{J}_k} \nabla_{\bx} F_j(\bx^k, \bb_j^k).
\end{align}
This formulation accommodates flexible stochastic gradient approximations, enabling the incorporation of advanced variance-reduction or momentum techniques, as detailed in \Cref{Sec: Algorithm implementations}.

Second, the auxiliary variable $\bB$ is updated via a vanilla gradient descent step \eqref{Eq: the update of B}; this suffices because the objective \eqref{Eq: simplified_lagrangian} with respect to $\bB$ is smooth with an inexpensive gradient, 
\begin{align}
   \nonumber \bB^{k+1} &= \bB^k - \tau \nabla_{\bB}\calL(\bx^{k+1}, \bB^k, \bLmbd^k)\\
    &= \bB^k - \tau \left( \nabla_{\bB} F(\bx^{k+1}, \bB^k) + \bLmbd^k \right). \label{Eq: the update of B}
\end{align}
Finally, the dual variable $\bLmbd$ is updated via proximal gradient ascent on $h^*$, using an extrapolation step introduced in \cite{chambolleFirstOrderPrimalDualAlgorithm2011}. The update solves the following optimization problem:
\begin{align}
    \bLmbd^{k+1} &= \text{prox}_{(1/\tau h^*)} \left( \bLmbd^k + \left(2\bB^{k+1} - \bB^k  \right) \right) \notag \\
    &=  \arg\min_{\bLmbd \in \mathcal{Y}^* } \left\{ h^*(\bLmbd) - \langle \bLmbd, 2\bB^{k+1} - \bB^k \rangle  + \frac{\tau}{2} \| \bLmbd - \bLmbd^k\|^2 \right\}. 
    \label{Eq: the update of lambda}
\end{align}

\subsection{Implementations}
\label{Sec: Algorithm implementations}
\paragraph{DualHash-StoM}
The heavy ball method \cite{Polyak1964} and Nesterov acceleration \cite{Nesterov1983AMF} has been extensively studied in large-scale nonconvex optimization \cite{gaoNonConvexStochasticComposite2024}.
Motivated by this, each iteration $k$,  we compute the stochastic gradient estimator (Line 3 of \Cref{alg:our_algorithm_sgdm}) as:
\begin{align}
\label{Eq: the stochastic approximate in SGDM}
    \bG^k = \nabla_{\bx} F(\bz^k, \bB^k; \calJ_k),\, \text{where } \bz^k = \bx^k + \beta_k (\bx^k - \bx^{k-1}).
\end{align} The update of \( \bx^k\) (Line 4 of Alg. \ref{alg:our_algorithm_sgdm}) is then:
\begin{align}
\label{Eq: the update of x in SGDM}
    \bx^{k+1} = \by^k - \eta_k \bG^k,\,\text{where } \by^k = \bx^k + \alpha_k (\bx^k - \bx^{k-1}). 
\end{align} The overall algorithm, DualHash-StoM, is detailed in \Cref{alg:our_algorithm_sgdm}.
\begin{remark}
The momentum acceleration framework in \eqref{Eq: the update of x in SGDM} covers these two accelerations: Nesterov acceleration corresponds to $\alpha_k = \beta_k \neq 0$; the heavy-ball method is obtained with $\beta_k = 0$; vanilla SGD is recovered when $\alpha_k = \beta_k = 0$. In this paper, we set $\alpha_k, \beta_k \in (0,1)$ for all $k \geq 1$. 
\end{remark}
\begin{algorithm}[t]
\caption{DualHash-StoM}
\label{alg:our_algorithm_sgdm}
\begin{algorithmic}[1]
\STATE \textbf{Initialize:} $\bx^1 = \bz^1 = \by^1 \in \calX $, \( \bB^1 \in \calB,\, \bLmbd^1\in \calY^*\), parameters $\{\eta_k\} $, $\tau$, $b_k$, $\{\alpha_k, \beta_k \in (0,1)\}$
\FOR{$k = 1, 2, \ldots$}               
  \STATE Compute \(\bG^k\) by \eqref{Eq: the stochastic approximate in SGDM}
  \STATE Update $\bx^{k+1}$ by \eqref{Eq: the update of x in SGDM}
  \STATE Update $\bB^{k+1}$ by \eqref{Eq: the update of B}
  \STATE Update $\bm{\Lambda}^{k+1}$ by \eqref{Eq: the update of lambda}
\ENDFOR
\STATE \textbf{Return:} $\bx_R$ where $R$ is uniformly sampled from $\{2, \ldots, T+1\}$
\end{algorithmic}
\end{algorithm}
\paragraph{DualHash-StoRM}
Alternatively, we employ the StoRM estimator \cite{cutkoskyMomentumBasedVarianceReduction2019} for the $\bx$-update. StoRM utilizes a recursive momentum mechanism to reduce gradient variance, which is motivated by SGDM. Unlike some other variance reduction techniques (e.g., SVRG \cite{johnsonAcceleratingStochasticGradient2013}, SARAH \cite{SARAH2017}) that require full-gradient or large-batch gradient computations, StoRM achieves variance reduction without such requirements. Given its practical effectiveness in deep learning \cite{cutkoskyMomentumBasedVarianceReduction2019,levy2021storm+}, we compute the stochastic gradient estimator (Line 3 of \Cref{alg:our_algorithm_storm}) as:
{\small
\begin{align}
\label{Eq: the stochastic approximate in STORM}
\bD^k =
\begin{cases} 
\nabla_{\bx} F(\bx^1, \bB^1; \mathcal{J}_1), &\text{if }  k = 1, \\
(1 - \rho_{k-1}) \big( \bD^{k-1} - \nabla_{\bx} F(\bx^{k-1}, \bB^{k-1};\mathcal{J}_{k-1})\big)+\nabla_{\bx} F(\bx^k, \bB^k; \mathcal{J}_k), & \text{if } k > 1.
\end{cases}
\end{align}}
The update of $\bx$ (Line 4 of \Cref{alg:our_algorithm_storm}) is then 
\begin{align}
\label{Eq: the update of x in STORM}
    \bx^{k+1} = \bx^k - \eta_k \bD^k.
\end{align} The overall algorithm, DualHash-StoRM, is detailed in \Cref{alg:our_algorithm_storm}.
\begin{algorithm}[t]
\caption{DualHash-StoRM}
\label{alg:our_algorithm_storm}
\begin{algorithmic}[1]
\STATE  \textbf{Initialize:} $\bx^1  \in \calX$, \( \bB^1 \in \calB,\, \bLmbd^1\in \calY^*\), parameters $\{\eta_k\} $, $\tau $, $b_k$, $\{\rho_k\in (0,1)\}$
\FOR{$k = 1, 2, \ldots$}
\STATE Compute $\bD^k$ by \eqref{Eq: the stochastic approximate in STORM}
\STATE Update $\bx^{k+1}$ by \eqref{Eq: the update of x in STORM}.
\STATE Update $\bB^{k+1}$ by \eqref{Eq: the update of B}.
\STATE Update $\bm{\Lambda}^{k+1}$ by \eqref{Eq: the update of lambda}.
\ENDFOR
\STATE \textbf{Return:} $\bx_R$ where $R$ is uniformly sampled from $\{2, \ldots, T+1\}$
\end{algorithmic}
\end{algorithm}

% \begin{remark}
% The update steps for $\bB$ and $\bm{\Lambda}$ imply the following optimality conditions, which may result in explicit closed-solution in deep hashing. The $\bB$-update in \eqref{Eq: the update of B} yields:
% \begin{align}
%     \label{Eq: the optimality condition of Bk}
%     \bB^{k} = \bB^{k-1} - \tau \left( \nabla_{\bB}F(\bx^{k}, \bB^{k-1}) + \bm{\Lambda}^{k-1} \right).
% \end{align}
% For the $\bm{\Lambda}$-update in \eqref{Eq: the update of lambda}, the proximal operator definition gives:
% \[
%     \tau \left( \bm{\Lambda}^{k} + \tau^{-1}(2\bB^{k} - \bB^{k-1}) - \bm{\Lambda}^{k} \right) \in \partial h^{*}(\bm{\Lambda}^{k}),
% \]
% which is equivalent to the existence of $\mathcal{G}_{\bm{\Lambda}}^{k} \in \partial h^{*}(\bm{\Lambda}^{k})$ such that:
% \begin{align}
%     \label{Eq: the optimality condition of lbda}
%     \mathcal{G}_{\bm{\Lambda}}^{k} = (2\bB^{k} - \bB^{k-1}) + \tau(\bm{\Lambda}^{k-1} - \bm{\Lambda}^{k}).
% \end{align}
% If the sequence \(\{ (\bx^k, \bB^k, \bm{\Lambda}^k) \}\) converges to \((\bx^*, \bB^*, \bm{\Lambda}^*)\), then by the outer semicontinuity of $\partial h^{*}(\cdot)$, we have $\bB^* \in \partial h^*(\bm{\Lambda}^*)$.
% This establishes \eqref{Eq: the sufficient and necessary condition for dual}, which is precisely the necessary condition for the second relation in \eqref{Eq: the sufficient and necessary condition}.
% \end{remark}